\documentclass[12pt,a4paper]{article}
\usepackage[utf8]{inputenc}
\usepackage{lmodern}
\usepackage[T1]{fontenc}
\usepackage[finnish]{babel}

\usepackage{graphicx}
\usepackage{epstopdf}

\begin{document}
\title{Tietokantojen perusteet, kevät 2020 \\
Harjoitustyön raportti}
\author{Ville Manninen \\014922187 \\Macville}
\date{\today}
\cleapage\maketitle
\begin{center}
\vspace{2cm}
\includegraphics{university-of-helsinki-2.eps} 
\end{center}
\thispagestyle{empty}

\newpage
\tableofcontents
\newpage
\section{Esittely}
Harjoitustyö on toteutettu käyttämällä NetBeans ympäristöä hyödyntäen SQlite tietokantaa sekä Java Database Connectivity (JDBC) ohjelmointirajapintaa. Harjoitustyöllä on mahdollista suorittaa seuraavat ominaisuudet.
\begin{itemize}
\item Tietokannan taulukoiden luominen.
\item Uusien asiakkaiden, paikkojen ja pakettien lisääminen.
\end{itemize}



\newpage
\section{Tietokantakaavio ja SQL-skeema}

\newpage
\section{Tehokkuustesti}
\subsection*{Ilman indeksejä}
\subsection*{Indeksien lisäämisen jälkeen}

\newpage
\section{Tietokannan eheys}
Tietokantaan ei ole mahdollista luoda samannimisiä asiakkaita tai paikkoja, eikä tietokantaan luoduilla paketeilla saa antaa samaa seurantakoodia. Samannimisiä syötteitä valvotaan SQliten toiminnolla UNIQUE, jonka avulla voidaan varmistaa että tietokantaan ei lisätä samannimisiä arvoja. Java sovellus varmistaa ettei sille annetut syötteet ole tyhjiä, joten tietokantaan ei tallenneta NULL arvoja.

Tietokantaan lisättäessä paketteja tai tapahtumia, jotka saavat arvokseen viiteavaimia, sovellus suorittaa kyselyn joka varmistaa että annetut viiteavaimet löytyvät tietokannasta.





\newpage
\section{Lähdekoodi}
\subsection*{GitHub}
\left \textbf{https://github.com/Viltska/sql-app-training} \right
\subsection*{Koodi}


\newpage

\end{document}